\section{Lord's Prayer, Communion, and Conclusion of Mass}

The Canon has offered the immaculate Victim and returned all glory to the Father through the Son. What follows is not an appendix to the sacrifice. The Church is taught to speak with the Son's filial voice; the consecrated species are broken and received; Communion is extended to the faithful when they approach; and the rite gathers its sacramental fruit before its several historical endings. These movements differ sharply in origin, voice, and theological weight, and must be read in their actual 1962 order.

\massmovement{\latin{Pater noster} and \latin{Libera nos}}{The Son's Prayer in the Sacrificing Church}{Fixed introduction, Lord's Prayer, response, and embolism}{\latin{Ordo Missae}, pp.~312--317; \latin{Ritus servandus} X.1--2}

The celebrant's \latin{Praeceptis salutaribus moniti, et divina institutione formati, audemus dicere} identifies the warrant for what follows before he begins the complete \latin{Pater noster}. In the 1962 arrangement he continues through \latin{Et ne nos inducas in tentationem}; the minister or choir answers \latin{Sed libera nos a malo}, and he supplies \latin{Amen} quietly. The dialogic ending is significant. The prayer belongs to the whole Church, yet the liturgical roles are not collapsed into an undifferentiated common recitation.

Matthew~6:5--15 places the prayer within Christ's correction of display, anxiety, and unforgiveness; Luke~11:1--13 joins the disciples' request to perseverance and the Father's gift. Its ``daily bread'' therefore retains the breadth of creaturely dependence even as Cyprian, Cyril, and Augustine receive it eucharistically (Cyprian, \emph{De dominica oratione} 18; Cyril, \emph{Mystagogical Catechesis} V.11--18; Augustine, Sermon~57.7). The movement from the Canon to the \latin{Pater} is not a descent from objective sacrifice into private devotion. Those gathered into Christ's oblation now pray in the form he gave, and the petition for forgiveness tests the peace with which they are about to approach one bread.

The \emph{Didache} securely witnesses thrice-daily Christian use (8.2--3). Optatus, \emph{De sacramentis}, and Augustine attest Latin eucharistic reception before Gregory the Great (Optatus, \emph{Contra Parmenianum} II.20, CSEL 26, p.~56; \emph{De sacramentis} V.4.24; Augustine, Sermon~227 and \emph{Epistle} 149.16). Gregory's October 598 letter documents and defends the Roman position of the Lord's Prayer immediately after the Eucharistic prayer; it does not establish that he first introduced the prayer into Eucharistic worship (\emph{Registrum} IX.26 in the modern numbering, MGH \emph{Epistolae} II, p.~59).

The celebrant then takes the paten and continues with \latin{Libera nos, quaesumus Domine, ab omnibus malis, praeteritis, praesentibus et futuris}. This Roman embolism is liturgical development, not an otherwise unrecorded saying of Christ. A precursor survives in the Verona collection (Muratori I, p.~359), while its saint list and gestures show layered transmission (Botte, \emph{Le canon}, p.~50, apparatus; Jungmann II, pp.~276--297). It prolongs the last petition of the Lord's Prayer into the communion of saints and the rite's immediate movement toward peace and fraction without becoming a second proper oration.

\massmovement{The Fraction of the Host}{One Host Broken, Christ Undivided}{Fixed fraction during the conclusion of \latin{Libera nos}}{\latin{Ordo Missae}, pp.~316--317; \latin{Ritus servandus} X.2}

Near the end of \latin{Libera nos}, the celebrant uncovers the chalice, genuflects, and breaks the Host over it. The larger parts are reunited on the paten while a small particle remains between his fingers. The action unfolds through the stable conclusion beginning \latin{Per eundem Dominum nostrum Iesum Christum, Filium tuum}, continuing with \latin{Qui tecum vivit et regnat in unitate Spiritus Sancti Deus}, and reaching the audible \latin{Per omnia saecula saeculorum} and its response \latin{Amen}. The break is thus enclosed within the prayer's address to the Father; it is not a dramatic representation staged apart from the oration.

Breaking bread names Eucharistic assembly in Acts~2:42 and 20:7, and Paul argues from one bread to the one ecclesial Body in 1~Corinthians~10:16--17. Those texts illuminate the received action without prescribing the later Roman division of the Host. Fraction was once a large practical act of distribution; the smaller 1962 gesture bears that ancestry after other practical functions had changed.

The sensible species are divided, not Christ. He remains wholly present under either species and under every part after division (Trent, Session XIII, ch.~3 and canon~3; \emph{ST} III, q.~76, a.~3). That doctrinal limit controls every later allegory of the broken Host. The fraction can disclose sacrificial giving and ecclesial unity only because the sacramental reality is not torn into pieces.

\massmovement{\latin{Pax Domini}}{Peace Proclaimed over the Chalice}{Fixed dialogue after fraction and before commingling}{\latin{Ordo Missae}, p.~317; \latin{Ritus servandus} X.2}

Holding the particle over the chalice, the celebrant signs from lip to lip three times while saying \latin{Pax Domini sit semper vobiscum}; the minister or choir responds \latin{Et cum spiritu tuo}. Only after this exchange does he place the particle in the chalice. The printed order matters: the dialogue of peace is neither the commingling formula nor the later transmission of the pax among ministers.

Christ gives peace in John~14:25--31 while going to the Father and confronting the ruler of this world. The Gospel setting prevents the dialogue from becoming a generic affirmation of concord. Here peace follows the Canon and fraction: it is announced under the signs of the sacrificed and living Lord before it is enacted, where appointed, in the distinct kiss of peace.

\massmovement{\latin{Haec commixtio} (Commingling)}{The Particle Enters the Chalice}{Fixed commingling immediately after the \latin{Pax Domini} response}{\latin{Ordo Missae}, p.~317; \latin{Ritus servandus} X.2}

After \latin{Et cum spiritu tuo}, the celebrant places the particle in the Precious Blood with the received formula \latin{Haec commixtio, et consecratio Corporis et Sanguinis Domini nostri Iesu Christi, fiat accipientibus nobis in vitam aeternam. Amen}. He then cleans and rejoins the consecrating fingers over the chalice, covers it, and genuflects. The ceremony is complete before the \latin{Agnus Dei} begins.

Earlier Roman uses of a particle, including the \latin{fermentum}, could manifest communion between celebrations; the later received gesture and formula have their own layered history. Neither history nor the word \latin{consecratio} licenses a second sacramental form. Commingling does not reconsecrate the species, make Christ whole again, or ritually restore life to him. Aquinas explicitly excludes such a reading while allowing the joined signs to signify the risen unity of Christ and the sacrament's ordered completion (\emph{ST} III, q.~83, a.~2, ad~2; a.~5). The rite joins what remains sacramentally distinct without implying that Christ had been divided.

\massmovement{\latin{Agnus Dei}, the Prayer for Peace, and the Pax}{The Lamb's Peace Received and Handed On}{Fixed \latin{Agnus Dei}; first peace prayer and transmission of the pax conditional, with Requiem variants}{\latin{Ordo Missae}, pp.~317--318; \latin{Ritus servandus} X.2--4}

With the chalice covered, the celebrant says the threefold \latin{Agnus Dei, qui tollis peccata mundi}, striking his breast at \latin{miserere nobis}, \latin{miserere nobis}, and \latin{dona nobis pacem}. Requiem Masses use \latin{dona eis requiem} twice and \latin{dona eis requiem sempiternam} the third time, and omit the breast striking. John the Baptist's testimony in John~1:29 identifies Jesus through sacrificial removal of sin; Revelation~5:6--14 places the slain Lamb within heavenly victory and worship. These contexts hold sacrifice, mercy, and peace together without reducing the chant to an interlude before reception.

The \emph{Liber pontificalis} retrospectively attributes the fraction chant to Sergius~I, and \emph{Ordo Romanus I}, no.~19, witnesses its Roman performance (Duchesne I, p.~376; PL 78:946). The third petition \latin{dona nobis pacem} is already present in tenth-century Saint-Martial, Winchester, and Reichenau tropers and becomes commoner in the eleventh century; no single composer is securely known (Blume--Bannister I, pp.~373, 385, nos.~385 and 419; Jungmann II, pp.~332--340).

The quiet prayer \latin{Domine Iesu Christe, qui dixisti Apostolis tuis: Pacem relinquo vobis} then turns from the celebrant's sins to the Church and the peace Christ promised. Where the pax is given, the celebrant kisses the altar and transmits it through the ordered ministry with \latin{Pax tecum} and \latin{Et cum spiritu tuo}. This is a distinct act from the earlier \latin{Pax Domini}. In Requiem Masses the prayer and transmission are omitted. The altar is the ritual point of departure: peace is received from Christ's sacrificial gift and handed on; it is not manufactured by the assembly's prior unanimity.

\massmovement{The Prayers before Communion}{Two Private Preparations, Two Different Accents}{Fixed quiet priestly prayers before Communion}{\latin{Ordo Missae}, pp.~318--319; \latin{Ritus servandus} X.3--4}

The priest remains bowed before the sacrament for \latin{Domine Iesu Christe, Fili Dei vivi} and \latin{Perceptio Corporis tui, Domine Iesu Christe}. Their juxtaposition should not disguise their distinct accents or histories. The first places approach within Christ's saving action and the priest's need for deliverance; the second confronts the possibility that reception may fail to bear its proper fruit. These are ministerial preparations in the singular voice, not substitutes for the Church's public Canon or for the faithful's own disposition.

\latin{Domine Iesu Christe, Fili Dei vivi} is attested in ninth-century Amiens and Le Mans books; \latin{Perceptio} occurs in the tenth-century Fulda tradition and the Ratoldus book (Leroquais, \emph{Eph\'em\'erides liturgiques} [1927], p.~444; \emph{Sacramentaires} I, p.~30; Richter--Sch\"onfelder, no.~24). Their separate witnesses disclose gradual medieval concentration on the priest's immediate sacramental accountability. Aquinas's distinction between Christ's objective presence and fruitful reception according to disposition supplies the doctrinal grammar without pretending to be the prayers' historical source (\emph{ST} III, q.~79).

\massmovement{The Communion of the Priest}{The Priest Receives under Both Species}{Fixed priestly Communion}{\latin{Ordo Missae}, pp.~319--321; \latin{Ritus servandus} X.4--5}

After genuflecting, the priest says \latin{Panem caelestem accipiam, et nomen Domini invocabo}. Holding the two larger parts above the paten, he strikes his breast and says three times \latin{Domine, non sum dignus, ut intres sub tectum meum: sed tantum dic verbo, et sanabitur anima mea}. He signs himself with the Host at \latin{Corpus Domini nostri Iesu Christi custodiat animam meam in vitam aeternam. Amen}, receives, joins his hands, and pauses. The pause after reception is one of the few places where the rubric itself makes room for unfilled attention rather than further speech.

Matthew~8:5--13 situates the centurion's words within Gentile faith in Christ's authority and a warning against presumptuous inheritance. The liturgical change from the servant of the Gospel narrative to \latin{anima mea} belongs to the received prayer's history, not to a claim about Matthew's wording. Longer Matthew~8 preparations appear from the tenth century; a \latin{puer} form is secure before 1174 at Modena, Roman and southern-Italian witnesses show further development, and the triple \latin{anima} form is standardized in 1570 (Jungmann II, pp.~355--359). The 1474 Missal still gives only one priestly recitation.

The priest begins \latin{Quid retribuam Domino pro omnibus quae retribuit mihi?}; while saying it, he uncovers the chalice, genuflects, rises, gathers any fragments from the corporal, and wipes the paten and his fingers over the chalice. Only then does he take the chalice and continue with \latin{Calicem salutaris accipiam, et nomen Domini invocabo. Laudans invocabo Dominum, et ab inimicis meis salvus ero}. He signs himself with the chalice at \latin{Sanguis Domini nostri Iesu Christi custodiat animam meam in vitam aeternam. Amen} and receives the Precious Blood with the commingled particle.

The priest's Communion is not a private possession attached to a public sacrifice. The one who sacramentally offers receives what has been offered; the rite moves from consecratory action into sacramental eating. A Mass at which no other person communicates remains a public act of Christ and the Church, while Trent also desires sacramental Communion by those present when properly disposed (Trent, Session XXII, ch.~6; \emph{Mediator Dei} 96--100). The two truths must not be played against one another.

\massmovement{The Communion of the Faithful}{The Sacrament Offered to the Faithful}{Conditional Communion of the faithful before the ablutions}{\latin{Ordo Missae}, pp.~321--322; \latin{Ritus servandus} X.6--9}

If the faithful communicate, the chalice is set safely within the corporal after the priest has received the Precious Blood. The celebrant prepares the consecrated particles, turns with one held above the ciborium or paten, and says \latin{Ecce Agnus Dei, ecce qui tollit peccata mundi}. He then says the complete \latin{Domine, non sum dignus} three times and distributes with \latin{Corpus Domini nostri Iesu Christi custodiat animam tuam in vitam aeternam. Amen}. The change from \latin{meam} at the priest's Communion to \latin{tuam} at distribution marks a genuine change of addressee.

The 1962 typical rubric at \latin{Ritus servandus} X.6 proceeds directly to \latin{Ecce Agnus Dei}; it does not print a second \latin{Confiteor}, \latin{Misereatur}, or \latin{Indulgentiam} in this sequence. The faithful's rite also has a history distinct from the priest's preparation. The earliest witness located in this audit for \latin{Ecce Agnus Dei} before Communion is the Synod of Aix in 1585, and the fuller complex is printed in the \emph{Rituale Romanum} of 1614 (Hardouin X, col.~1525; Jungmann II, pp.~367--374). It must not be projected backward as an unchanged apostolic or early-medieval unit.

Distribution itself is ancient. Justin describes the Eucharist being given and carried to absent members (\emph{First Apology} 65, 67). Cyril joins the communicant's \latin{Amen}, adoration, and care for every particle, but his Jerusalem gestures are evidence for fourth-century mystagogy, not 1962 Roman rubrics (\emph{Mystagogical Catechesis} V.19--23). Augustine's command to become what the communicant sees connects the sacramental Body received with the ecclesial body being formed (Sermons 227 and 272). First Corinthians~10--11 supplies the severe canonical frame: real communion with Christ judges faction, contempt, and undiscerning reception. Conditional occurrence therefore does not make the faithful's Communion theologically peripheral.

\massmovement{The Purification and Ablutions}{Reverence after Reception}{Fixed purification after the last Communion}{\latin{Ordo Missae}, p.~322; \latin{Ritus servandus} X.5, 7}

After the last person has communicated, the celebrant returns any remaining particles, gathers fragments, and only then begins the ablutions. After priestly Communion without distribution, the same stable sequence follows immediately. He says \latin{Quod ore sumpsimus, Domine, pura mente capiamus: et de munere temporali fiat nobis remedium sempiternum} as wine is poured into the chalice. Wine and water then wash the joined fingers over the chalice while he says \latin{Corpus tuum, Domine, quod sumpsi, et Sanguis, quem potavi, adhaereat visceribus meis}. He consumes the ablutions, dries and arrays the vessels, and returns the veiled chalice to the altar.

John~6:12 can illuminate the care of fragments by analogy, but the gathering of multiplied loaves is not the historical institution of chalice ablutions. Cyril's warning against losing a particle is a direct early witness to Eucharistic reverence (\emph{Mystagogical Catechesis} V.21--22). Such care follows the Church's confession that Christ remains present while the species remain; it is not a detachable scruple or an attempt to repeat Communion.

The prayers' ancestry differs from their received private placement. \latin{Quod ore} occurs as a Postcommunion in the Verona collection (Muratori I, p.~366), and \latin{Corpus tuum} has an ancient public-prayer ancestor in the Gothic Missal (Muratori II, p.~653). Explicit Western cleansing appears in ninth- and tenth-century witnesses, wine in the eleventh century, and the familiar wine-and-water pattern in the Dominican Ordinarium of 1256 before Roman standardization in 1570 (Jungmann II, pp.~392--425). The physical action and the prayers converged gradually; neither may be assigned wholesale to apostolic practice.

\massmovement{The Communion Antiphon, Postcommunion, and \texorpdfstring{\mbox{\latin{Oratio super populum}}}{Oratio super populum}}{The Day Interprets the Communion It Has Received}{Proper and conditional texts within a fixed frame}{\latin{Ordo Missae}, pp.~322--323; \latin{Ritus servandus} XI}

Once purified, the celebrant reads the appointed Communion antiphon. He returns to the center, kisses the altar, and exchanges \latin{Dominus vobiscum} and \latin{Et cum spiritu tuo}; at the Missal, \latin{Oremus} opens the appointed Postcommunion prayer or prayers, to which the response is \latin{Amen}. In sung Mass the antiphon belongs to the choir's Communion action even while the complete-Missal presentation assigns its recitation to the celebrant. \emph{Ordo Romanus I}, nos.~20--21, already witnesses Communion psalmody followed by prayer after Communion (Andrieu II, pp.~103--107; PL 78:947--948A).

The antiphon and Postcommunion are stable positions, not invariant texts. Their own biblical, hagiographic, or ecclesial sources must govern their interpretation in each Mass; no universal paraphrase can predetermine the particular fruit for which the proper oration asks. On the appointed Lenten ferias, the prayer over the people intervenes after the Postcommunion and before the final greeting (\latin{Ritus servandus} XI.2). This is a conditional part of the received conclusion, not another Postcommunion.

The order from antiphon to oration prevents sacramental fruit from being reduced either to private feeling or to automatic moral transformation. The gift is objective; its ecclesial and personal fruit still calls for thanksgiving, disposition, and divine action (Trent, Session XIII, chs.~7--8; \emph{ST} III, q.~79, aa.~1, 3--8). Only after this proper interpretation of Communion does the Mass turn toward its historically layered endings.

\massmovement{\latin{Ite, missa est}, \latin{Benedicamus Domino}, or \latin{Requiescant in pace}}{The Celebration Is Released}{Conditional forms within a fixed final dialogue}{\latin{Ordo Missae}, p.~323; \latin{Ritus servandus} XII.1; General Rubric 507}

The celebrant returns to the center, kisses the altar, and exchanges the final \latin{Dominus vobiscum} and \latin{Et cum spiritu tuo}. The ordinary form is \latin{Ite, missa est} with \latin{Deo gratias}; in temporal Masses within the Easter octave, both receive a double \latin{Alleluia}. General Rubric 507a appoints \latin{Benedicamus Domino} when a procession follows; Maundy Thursday also uses it before the solemn reposition. Requiem Mass uses \latin{Requiescant in pace} with the response \latin{Amen}. These are true rubrical alternatives, not ornamental variations on one English slogan.

\latin{Missa} functions as dismissal in \emph{Ordo Romanus I}, no.~21 (Andrieu II, pp.~107--108; PL 78:948). Earlier uses such as \emph{Peregrinatio Aetheriae} 25.1--2 likewise concern departure, not the later name of the complete Eucharistic rite. Historical and linguistic discipline therefore forbids deriving an entire theology of mission from the syllables of \latin{missa}.

\textit{Theological synthesis, not lexical derivation.} The dismissal can nevertheless be read within the risen Christ's sending of witnesses in Matthew~28:18--20 and Luke~24:44--49. Those biblical commissions illuminate what sacramental Communion demands of a blessed Church; they do not supply the etymology of \latin{Ite, missa est} or prove that its first users encoded a modern mission program in the formula. Mission is the fruit of the whole action just completed, not a pun imposed on its final words.

\massmovement{\latin{Placeat tibi, sancta Trinitas}}{The Priest's Service Submitted to Judgment}{Fixed quiet \latin{Placeat} after the dismissal}{\latin{Ordo Missae}, pp.~325--326; \latin{Ritus servandus} XII.1}

Facing the altar and deeply bowed, the priest says \latin{Placeat tibi, sancta Trinitas, obsequium servitutis meae}. This private conclusion turns away from any assumption that exact performance certifies its own acceptance. The Mass's sacramental efficacy rests on Christ, while the minister still submits his service to divine mercy. That distinction preserves both objective sacramental action and personal accountability without making the prayer a second offertory or Canon.

Exact \latin{Placeat} witnesses occur in ninth-century Amiens and Le Mans books, followed by Fulda, Ratoldus, and Bernold (Jungmann II, pp.~437--438). The prayer is therefore a historically layered private ending, not an apostolic formula and not the source of the dismissal that precedes it. Its position after the public exchange makes the contrast audible: the appointed response has sounded; the priest remains bowed before the altar.

\massmovement{\latin{Benedictio in fine Missae} (Blessing)}{Beneath the Returning Sign of the Cross}{Normally fixed, with the omissions governed by General Rubric 508}{\latin{Ordo Missae}, p.~326; \latin{Ritus servandus} XII.1; General Rubric 508}

After \latin{Placeat}, the priest kisses the altar, raises and joins his hands, and begins \latin{Benedicat vos omnipotens Deus}. Turning to the people, he blesses them with \latin{Pater, et Filius, et Spiritus Sanctus}; the response is \latin{Amen}. General Rubric 508 omits the blessing after \latin{Benedicamus Domino} or \latin{Requiescant in pace}. Its normal presence must therefore not be made a condition for the validity or integrity of the Eucharistic sacrifice.

A final presbyteral blessing is historically later than the ancient Roman dismissal. A late-eleventh-century Brescia witness, an Albi synodal formula in 1230, and the local Dominican order of 1256 precede the Roman sequence fixed by Pius~V in 1570 and refined by Clement~VIII in 1604 (Jungmann II, pp.~439--446). The biblical priestly blessing in Numbers~6:22--27 and apostolic conclusions such as 2~Corinthians~13:13 illuminate the act theologically; they do not establish the date or wording of this Roman formula.

The opening Trinitarian sign now returns as an act bestowed on the people. This is structural synthesis: the worshippers depart under God's Name and action, not merely with a recollection of what they themselves have done.

\massmovement{\latin{Ultimum Evangelium} (Last Gospel)}{Incarnation after Communion}{Normally John~1:1--14, with proper substitutions and omissions}{\latin{Ordo Missae}, pp.~326--328; \latin{Ritus servandus} XII.2--8; General Rubrics 509--510}

At the Gospel side the celebrant again exchanges \latin{Dominus vobiscum} and \latin{Et cum spiritu tuo}. While announcing \latin{Initium sancti Evangelii secundum Ioannem}, he signs the altar or book and then his forehead, lips, and breast; the response \latin{Gloria tibi, Domine} follows, and the normal text begins \latin{In principio erat Verbum}. He genuflects at \latin{Et Verbum caro factum est}, and the response at the end is \latin{Deo gratias}. General Rubric 509 appoints a proper Last Gospel for Palm Sunday Masses not following the blessing and procession of branches. General Rubric 510 omits the Last Gospel whenever \latin{Benedicamus Domino} was said under General Rubric 507a---including the evening Mass of the Lord's Supper before solemn reposition---and also at the third Mass of Christmas, the Palm Sunday procession Mass, the Easter Vigil, the specified Requiem Masses followed by absolution over the \latin{tumulum}, and the specified consecration Masses.

John~1:1--14 is not a closing motto but the opening movement of a complete Gospel: the eternal Word is God, creates and illuminates, is rejected and received, and truly becomes flesh. Augustine's exposition guards both immutable divinity and assumed humanity (\emph{Tractates on John} 1.8--15; 2.2--3). The genuflection answers the Incarnation confessed earlier in the Creed without suggesting a new Incarnation or a second Gospel liturgy at Mass's end.

Its position is historically layered. The Synod of Seligenstadt in 1022 witnesses devotional use of the Prologue but not an end-of-Mass Gospel. The earliest conclusion-of-Mass witness located in this audit is the Dominican Ordinarium of 1256; Li\`ege about 1285 and Durandus show further spread. The approved Lyon Missal of 1550 and Jesuit General Constitutions of 1558 still lack it, and the 1474 Roman Missal ends after \latin{Placeat} and blessing (Lippe, HBS 17, pp.~210--212). Pius~V standardized an inherited but nonuniversal practice in 1570; he did not compose the Prologue or invent its devotional use (Jungmann II, pp.~447--451).

As the final received lens, the Prologue places every preceding sacramental sign under the reality of the Word made flesh. The Gospel first approached with procession and incense is now borne outward in the life of a Church that has heard, offered, and received. This is a theological relation within the whole Mass, not a claim that John composed his Prologue for this Roman position. The 1962 \latin{Ordo Missae} ends here, ordinarily with \latin{Deo gratias}; the externally prescribed Leonine prayers treated later remain outside the Mass.
